\title{MetaICL: Learning to Learn In Context}

\begin{abstract}
We introduce MetaICL\ (\textbf{Meta}-training for \textbf{I}n-\textbf{C}ontext \textbf{L}earning), a new meta-training framework for few-shot learning where a pretrained language model is tuned to do in-context learning on a large set of training tasks. This meta-training enables the model to more effectively learn a new task in context at test time, by simply conditioning on a few training examples with no parameter updates or task-specific templates.

We experiment on a large, diverse collection of tasks consisting of 142 NLP datasets including classification, question answering, natural language inference, paraphrase detection and more, across seven different meta-training/target splits.

MetaICL\ outperforms a range of baselines including in-context learning without meta-training and multi-task learning followed by zero-shot transfer. We find that the gains are particularly significant for target tasks that have domain shifts from the meta-training tasks, and that using a diverse set of the meta-training tasks is key to improvements. We also show that MetaICL\ approaches (and sometimes beats) the performance of models fully finetuned on the target task, and outperforms much bigger models with nearly 8x parameters. Finally, we show that MetaICL\ is complementary to human-written instructions, and the best performance can be achieved by combining both approaches.
\end{abstract}

\section{Introduction}\label{sec:intro}

Large language models (LMs) have recently been shown to be able to do {\em in-context learning}~\citep{brown2020language}, where they learn a new task simply by conditioning on a few training examples and predicting which tokens best complete a test input. This type of learning is attractive because the model learns a new task through inference alone, without any parameter updates. However, performance significantly lags behind supervised finetuning, results are often high variance~\citep{zhao2021calibrate,perez2021true}, and it can be difficult to engineer the templates that convert existing tasks to this format. 

In this paper, we address these challenges by introducing MetaICL: \textbf{Meta}-training for \textbf{I}n-\textbf{C}ontext \textbf{L}earning. MetaICL\ tunes a pretrained language model on a large set of tasks to learn how to in-context learn, and is evaluated on strictly new unseen tasks. Each meta-training example matches the test setup---it includes $k+1$ training examples from one task that will be presented together as a single sequence to the language model, and the output of the final example is used to calculate the cross-entropy training loss.

Simply finetuning the model in this data setup directly leads to better in-context learning---the model learns to recover the semantics of the task from the given examples, as must be done for in-context learning of a new task at test time. This approach is related to recent work that uses multi-task learning for better zero-shot performance at test time~\citep{ khashabi2020unifiedqa,zhong2021adapting,mishra2022cross,wei2022finetuned,sanh2022multitask}. However, MetaICL\ is distinct as it allows learning new tasks from $k$ examples alone, without relying on a task reformatting (e.g., reducing everything to question answering) or task-specific templates (e.g., converting different tasks to a language modeling problem).

We experiment on a large, diverse collection of tasks taken from \citet{ye2021crossfit} and \citet{khashabi2020unifiedqa}, including 142 text classification, question answering, natural language inference and paraphrase detection datasets. We report seven different settings, all with no overlap between meta-training and target tasks. This leads to 52 unique target tasks in total, which is the largest among all recent related work to the best of our knowledge.

Experimental results show that MetaICL\ consistently outperforms baselines including (1) a variety of LM in-context learning baselines without meta-training~\citep{brown2020language,zhao2021calibrate,holtzman2021surface,min2022noisy}, and (2) multi-task learning followed by zero-shot transfer~\citep{zhong2021adapting,wei2022finetuned,sanh2022multitask}. Gains over multi-task zero-shot transfer are particularly significant when meta-training tasks and target tasks are dissimilar, e.g. there are large differences in task formats, domains, or required skills. This demonstrates that MetaICL\ enables the model to recover the semantics of the task in context during inference even when the target does not share similarities with meta-training tasks.

MetaICL\ often gets close to (and sometimes beats) the performance of models trained with supervised finetuning on the target datasets, and perform as well as models with 8x parameters. We also perform extensive ablations to identify key ingredients for success of MetaICL\ such as the number and diversity of meta-training tasks. Finally, we demonstrate MetaICL\ without any templates is better than recent work using human-written natural instructions, while the best performance is achieved by combining both approaches. Code and data are publicly released at 
    \href{https://github.com/facebookresearch/MetaICL}{\nolinkurl{github.com/facebookresearch/MetaICL}}
.

\section{Related Work}\label{sec:related}

\paragraph{In-context learning}

\citet{brown2020language} propose to use a language model (LM) conditioned on a concatenation of training examples for few-shot learning with no parameter updates. It has been further improved by later work~\citep{zhao2021calibrate, holtzman2021surface,min2022noisy}, showing promising results on a variety of tasks. However, in-context learning with an LM achieves poor performance when the target task is very different from language modeling in nature or the LM is not large enough. Moreover, it can have high variance and poor worst-case accuracy~\citep{perez2021true,lu2021fantastically}.

Our paper is based on the core idea of in-context learning by conditioning on training examples. We show that, by explicitly training on an in-context learning objective, MetaICL\ achieves substantial improvements even with smaller LMs. 

\paragraph{Meta-training via multi-task learning} Our work is broadly inspired by a large body of work in meta-learning~\citep{vilalta2002perspective,finn2017model} and multi-task learning~\citep{evgeniou2004regularized,ruder2017overview}. Prior work has shown that multi-task learning on a large collection of tasks leads to better performance on a new task, either when tested zero-shot~\citep{khashabi2020unifiedqa,zhong2021adapting,mishra2022cross,wei2022finetuned} or when further finetuned~\citep{aghajanyan2021muppet,ye2021crossfit}. In particular, the former is closely related to our work, as it eliminates the need for parameter updates on a target task. However, these zero-shot models are either limited to tasks sharing the same format as training tasks (e.g., a question answering format)~\citep{khashabi2020unifiedqa,zhong2021adapting}, or rely heavily on task-specific templates~\citep{mishra2022cross,wei2022finetuned,sanh2022multitask}

which are difficult to engineer due to high variance in performance from very small changes~\citep{mishra2021reframing}.

In this paper, we propose a meta-training method for better in-context learning that improves few-shot performance. We show that it effectively learns semantics of a new task with no manual effort, significantly outperforming zero-shot transfer methods.\footnote{We show that MetaICL\ without instructions is still better than zero-shot transfer with instructions, but by using instructions, performance of MetaICL\ further improves (Section~\ref{subsec:ablations}).}
Furthermore, while \citet{wei2022finetuned} show that meta-training helps only when the model has 68B or more parameters, our experiments demonstrate improvements with a much smaller model (770M).

\citet{chen2022meta}, concurrently to our work, propose meta-training for in-context learning. Our approach differs in a number of ways:

we remove requirements of human-written templates or instructions, and include more diverse tasks, stronger baselines, and extensive experiments in much larger scale with many meta-training/target splits.

\section{MetaICL}\label{sec:method}

We introduce MetaICL: \textbf{Meta}-training for \textbf{I}n-\textbf{C}ontext \textbf{L}earning.
Table~\ref{tab:overview} provides an overview of the approach. The key idea is to use a multi-task learning scheme over a large collection of meta-training tasks, in order for the model to learn how to condition on a small set of training examples, recover the {\em semantics} of a task, and predict the output based on it. Following previous literature~\citep{brown2020language}, the training examples are concatenated and provided as an single input to the model, which is feasible for $k$-shot learning (e.g., $k=16$). At test time, the model is evaluated on an unseen target task that comes with $k$ training examples, and inference directly follows the same data format as in meta-training.

\subsection{Meta-training}\label{subsec:meta-training}

The model is meta-trained on a collection of tasks which we call meta-training tasks. For every iteration, one meta-training task is sampled, and $k+1$ training examples $(x_1,y_1), \cdots, (x_{k+1},y_{k+1})$ are sampled from the training examples of the chosen task. We then supervise the model by feeding the concatenation of $x_1, y_1, \cdots, x_k, y_k, x_{k+1}$ to the model as an input and train the model to generate $y_{k+1}$ using a negative log likelihood objective. This simulates in-context learning at inference where the first $k$ examples serve as training examples and the last $(k+1)$-th example is regarded as the test example.

\subsection{Inference}\label{subsec:inference}

For a new target task, the model is given $k$ training examples $(x_1, y_1), \cdots, (x_k, y_k)$ as well as a test input $x$. It is also given a set of candidates $\mathcal{C}$ which is either a set of labels (in classification) or answer options (in question answering). As in meta-training, the model takes a concatenation of $x_1, y_1, \cdots, x_k, y_k, x$ as the input, and compute the conditional probability of each label $c_i \in \mathcal{C}$.
The label with the maximum conditional probability is returned as a prediction.

\subsection{Channel MetaICL}\label{subsec:channel-mic}
We introduce a noisy channel variant of MetaICL\ called Channel MetaICL, following \citet{min2022noisy}.
In the noisy channel model, $P(y|x)$ is reparameterized to
$\frac{P(x|y)P(y)}{P(x)} \propto P(x|y)P(y)$. We follow \citet{min2022noisy} in using $P(y)=\frac{1}{|\mathcal{C}|}$ and modeling $P(x|y)$ which allows us to use the channel approach by simply flipping $x_i$ and $y_i$. Specifically, at meta-training time, the model is given a concatenation of $y_1, x_1, \cdots, y_k, x_k, y_{k+1}$ and is trained to generate $x_{k+1}$. At inference, the model computes $\mathrm{argmax}_{c \in \mathcal{C}} P(x|y_1,x_1,\cdots,y_k,x_k,c)$. 

\section{Experimental Setup}\label{sec:exp-setup}

\subsection{Datasets}\label{subsec:dataset}
We use a large collection of tasks taken from \textsc{CrossFit}~\citep{ye2021crossfit} and \textsc{UnifiedQA}~\citep{khashabi2020unifiedqa}.
We have 142 unique tasks in total, covering a variety of problems including text classification, question answering (QA), natural language inference (NLI) and paraphrase detection. All tasks are in English.

We experiment with seven distinct settings as shown in Table~\ref{tab:data-summary}, where there is no overlap between the meta-training and target tasks. The number of unique target tasks in total is 52, which is significantly larger than other relevant work~\citep{khashabi2020unifiedqa,zhong2021adapting,mishra2022cross,wei2022finetuned,sanh2022multitask}. Each target task is either classification or multi-choice, where a set of candidate options ($\mathcal{C}$ in Table~\ref{tab:overview}) is given.

 \textbf{HR$\rightarrow$LR} (High resource to low resource): We experiment with a setting where datasets with 10,000 or more training examples are used as meta-training tasks and the rest are used as target tasks. We think using high resource datasets for meta-training and low resource datasets as targets is a realistic and practical setting for few-shot learning.

 \textbf{X$\rightarrow$X (X=\{Classification, QA\})}: We experiment with two settings with meta-training and target tasks sharing the task format, although with no overlap in tasks.

 \textbf{Non-X$\rightarrow$X (X=\{Classification, QA, NLI, Paraphase\})}: Lastly, we experiment with four settings where meta-training tasks do not overlap with target tasks in task format and required capabilities. These settings require the most challenging generalization capacities.

Each setting has a subset of target tasks with no domain overlap with any meta-training tasks (e.g., finance, poem, climate or  medical). We report both on all target tasks or on target tasks with no domain overlap only. Full details of the settings and datasets with citations are provided in Appendix~\ref{app:dataset}.

\subsection{Baselines}\label{subsec:baselines}
We compare MetaICL\ and Channel MetaICL\ with a range of baselines, as summarized in Table~\ref{tab:methods}.

 \textbf{0-shot}: We use a pretrained LM as it is and run zero-shot inference, following \citet{brown2020language}.

 \textbf{In-context}: We use the pretrained LM as it is and use in-context learning by conditioning on a concatenation of $k$ training examples, following \citet{brown2020language}.

 \textbf{PMI 0-shot, PMI In-context}: We use the PMI method from \citet{holtzman2021surface,zhao2021calibrate} for 0-shot and In-context learning.

 \textbf{Channel 0-shot, Channel In-context}: We use the noisy channel model from \citet{min2022noisy} for 0-shot and In-context learning.

 \textbf{Multi-task 0-shot}: We train the LM on the same meta-training tasks without in-context learning objective, i.e., maximize $P(y|x)$ without $k$ other training examples, and then use zero-shot transfer on a target task. This is equivalent to MetaICL\ with $k=0$. This is a typical multi-task learning approach from previous work~\citep{khashabi2020unifiedqa,zhong2021adapting,wei2022finetuned}.

 \textbf{Channel Multi-task 0-shot}: We have a channel variant of Multi-task 0-shot.

 \textbf{Fine-tune}: We fine-tune the LM on an individual target task. This is not directly comparable to other methods as parameter updates are required for every target task.

 \textbf{Fine-tune w/ meta-train}: We train the LM on meta-training tasks first and then further fine-tuned it on a target task. This is not directly comparable to other methods for the same reason as above.

\newcommand{{\protect\color{green} [P]}}{{\protect\color{green} [P]}}
\newcommand{{\protect\color{maroon} [H]}}{{\protect\color{maroon} [H]}}

\newcolumntype{C}[1]{>{\PreserveBackslash\centering}p{#1}}

\subsection{Evaluation}\label{subsec:evaluation}
We use Macro-F1\footnote{More suitable than accuracy for imbalanced classification.} and Accuracy as evaluation metrics for classification tasks and non-classification tasks, respectively. 

For a target task, we use $k=16$ training examples, sampled uniformly at random. We relax the assumption of perfect balance between labels on $k$ training examples, following \citet{min2022noisy}.
Because in-context learning is known to have high variance~\citep{zhao2021calibrate,perez2021true,lu2021fantastically}, we use 5 different sets of $k$ training examples. We first compute the average and the worst-case performance over seeds for every target task, and then report the macro-average of them over all target tasks.

\subsection{Experiment Details}\label{subsec:impl-details}
As a base LM, we use GPT-2 Large ~\citep{radford2019language} which consists of 770M parameters.\footnote{Appendix~\ref{app:abl_lm_size} reports performance for other LM sizes.} For baselines without meta-training (raw LMs), we also compare with GPT-J~\citep{wang2021gpt}, which is the largest public causal LM at the time of writing, consisting of 6B parameters.

\paragraph{Elimination of templates} Prior work uses human-authored templates to transform the input-output pair to a natural language sentence~\citep{zhong2021adapting,mishra2022cross,wei2022finetuned,chen2022meta}. They require expensive manual effort (as 136 different templates are required for 136 tasks in this paper) and cause unstable model performance due to many different ways of writing~\citep{mishra2021reframing}.

We eliminate templates, using the given input (or a concatenation of inputs if there are multiple) and label words provided in the original datasets.\footnote{In our preliminary experiments, we explored templates taken from prior work, but found that they do not consistently improve few-shot performance, even when they do improve zero-shot performance.}

A comparison of input-output schemes from prior work and our approach is shown in Table~\ref{tab:templates}.

\paragraph{Training details} All implementation is done in PyTorch~\citep{paszke2019pytorch} and Transformers~\citep{wolf-etal-2020-transformers}. For meta-training, we use up to 16,384 training examples per task. We use a batch size of $8$, learning rate of $1 \times 10^{-5}$ and a sequence length of $1024$. For multi-task 0-shot baselines (the baselines with no in-context learning), we use a sequence length of $256$. We train the model for $30,000$ steps.\footnote{We also explored training longer, but it did not improve performance.} To save memory during meta-training, we use an 8-bit approximation~\citep{dettmers20228} of an Adam optimizer~\citep{kingma2015adam} and mixed precision~\citep{micikevicius2017mixed}. Training was done for 4.5 hours with eight 32GB GPUs. This is drastically more efficient than recent prior work, e.g., 270 hours of a 512GB TPU in \citet{sanh2022multitask}.

More details about preprocessing and training can be found in Appendix~\ref{app:impl-details}.

\section{Experimental Results}\label{sec:exp-result}

\subsection{Main Results}\label{subsec:main-results}

Table~\ref{tab:main-result} reports the full results using GPT-2 Large, where we compute the average and the worst-case performance of every target task and report the macro-average over them. The top and the bottom respectively evaluate on all target tasks and target tasks in unseen domains only.

\paragraph{Our baselines are strong} We first discuss the results of ours baselines. Among raw LMs without meta-training (the first six rows of Table~\ref{tab:main-result}), we observe that channel in-context baselines are the most competitive, consistent with findings from \citet{min2022noisy}.
We then find that Multi-task 0-shot baselines do not outperform the best raw LM baseline in most settings, despite being supervised on a large set of meta-training tasks. This somewhat contradicts findings from \citet{wei2022finetuned,sanh2022multitask}.
This is likely for two reasons. First, our models are much smaller than theirs (770M vs. 11B--137B); in fact, \citet{wei2022finetuned} reports Multi-task 0-shot starts to be better than raw LMs only when the model size is 68B or larger. Second, we compare with much stronger channel baselines which they did not;

Multi-task 0-shot outperforms non-channel LM baselines but not channel LM baselines.

\paragraph{MetaICL\ outperforms baselines}
MetaICL\ and Channel MetaICL\ consistently outperform a range of strong baselines.

In particular, Channel MetaICL\ achieves the best performance in 6 out of 7 settings. Gains are particularly significant in the HR$\rightarrow$LR, non-NLI$\rightarrow$NLI and non-Para$\rightarrow$Para settings (6--15\% absolute). This is noteworthy because HR$\rightarrow$LR\ targets the common low-resource case where new tasks have very few labeled examples, and the other two represent large data distribution shifts where the test tasks are relatively different from the meta-training tasks. This demonstrates that MetaICL\ can infer the semantics of new tasks in context even when there are no closely related training tasks.

While MetaICL\ significantly outperforms baselines in most settings, it only marginally outperforms Multi-task 0-shot in the QA$\rightarrow$QA setting, as an exception. This is likely because the meta-training and target tasks are relatively similar, allowing the Multi-task 0-shot baseline to achieve very strong performance. Nonetheless, performance of Multi-task 0-shot in QA significantly drops when the model is trained on non-QA tasks, while performance of MetaICL\ drops substantially less.

\paragraph{Gains are larger on unseen domains} Gains over Multi-task 0-shot are more significant on target tasks in unseen domains. In particular, Multi-task 0-shot is generally less competitive compared to raw LM baselines, likely because they require more challenging generalization.
MetaICL\ suffers less from this problem and is consistently better or comparable to raw LM baselines across all settings.

\paragraph{Comparison to fine-tuning}
MetaICL\ matches or sometimes even outperforms fine-tuned models without meta-training.

This is a promising signal, given that no prior work has shown models with no parameter updates on the target can match or outperform supervised models. Nonetheless, fine-tuning with meta-training exceeds both MetaICL\ and fine-tuning without meta-training, because meta-training helps in supervised learning as it does in in-context learning.

This indicates that there is still room for improvement in methods that allow learning without parameter updates .

\paragraph{Comparison to GPT-J} In Table~\ref{tab:comparison-gptj}, we compare GPT-2 Large based models with raw LM baselines based on GPT-J which consists of 6B parameters. MetaICL, despite being 8x smaller, outperforms or matches GPT-J baselines. 

\subsection{Ablations}\label{subsec:ablations}

\paragraph{Varying number of training examples} We vary the number of training examples ($k$) from 0, 4, 8, 16 to 32. In-context learning with $k=0$ is equivalent to the zero-shot method. Results are shown in Figure~\ref{fig:ablate_k}. Increasing $k$ generally helps across all models, and Channel MetaICL outperforms the raw in-context learning over all values of $k$. We additionally find that the performance tends to saturate when $k$ is closer to $16$, likely because the sequence length limit of the language model makes it hard to encode many training examples.

\paragraph{Number of meta-training tasks} To see the impact of the number of meta-training tasks, we subsample $\{7,15,30\}$ meta-training tasks out of 61 in the HR$\rightarrow$LR\ setting. For each, we use ten different random seeds to additionally see the impact of the choice of meta-training tasks.

Figure~\ref{fig:ablate_size} reports the results. On average, performance generally increases as the number of tasks increase, which is consistent with results in \citet{mishra2022cross,wei2022finetuned}.
Across different numbers of meta-training tasks, Channel MetaICL\ consistently outperforms other models. Nonetheless, there is nonnegligible variance across different choices of meta-training (the bottom of Figure~\ref{fig:ablate_size}), indicating that a choice of meta-training gives substantial impact in performance.

\paragraph{Diversity in meta-training tasks} We hypothesize that the diversity in meta-training tasks may impact performance of MetaICL. To verify this hypothesis, we create two settings by subsampling 13 out of 61 meta-training datasets in the HR$\rightarrow$LR\ setting. One setting is diverse in their task formats and required capacities: QA, NLI, relation extraction, sentiment analysis, topic classification, hate speech detection and more. The other setting is less diverse, including tasks related to sentiment analysis, topic classification and hate speech detection only. A full list of datasets is reported in Appendix~\ref{app:dataset}. Using these two settings, we compare multi-task zero-shot transfer baselines and MetaICL.

Results are reported in Table~\ref{tab:ablate_diversity}.
We find that MetaICL\ with a diverse set outperforms MetaICL\ with a non-diverse set by a substantial margin. This shows that diversity among meta-training tasks is one of substantial factors for the success of MetaICL.

In Appendix~\ref{app:which-tasks-helpful}, we include ablations that provide more insights on the choice of meta-training tasks, such as (1) high quality data with diverse domains tend to help (e.g., GLUE family~\citep{wang2018glue}) and (2) adversarially collected data tends to be unhelpful. However, more systematic studies on how to choose the best meta-training tasks and how they relate to particular target tasks should be done, which we leave for future work.

\paragraph{Are instructions necessary?} Most recent work has used human-written natural instructions for zero- or few-shot learning~\citep{mishra2022cross, wei2022finetuned,sanh2022multitask}. While we argue for not using instructions to avoid manual engineering and high variance, we also ask: {\em are instructions still useful with MetaICL?} On one hand, learning to condition on $k$ examples may remove the necessity of instructions. On the other hand, instructions may still be complementary and provide the model with extra useful infomration.

We aim to answer this question by using 32 meta-training tasks and 12 target tasks from the HR$\rightarrow$LR\ setting for which human-written instructions are available in \citet{sanh2022multitask}.\footnote{\url{github.com/bigscience-workshop/promptsource}}
We have two variants: (a) using one instruction per meta-training task, and (b) using all available instructions including 267 instructions in total (8.3 per meta-training task) which \citet{sanh2022multitask} found to be better than (a). We then compare MetaICL\ and a range of baselines with and without instructions.

Results are reported Table~\ref{tab:ablate_inst}.
As in \citet{wei2022finetuned} and \citet{sanh2022multitask}, Multi-task 0-shot outperforms the raw-LM 0-shot baseline. However, MetaICL\ with no instructions is better than Multi-task 0-shot with instructions. Furthermore, MetaICL\ achieves further improvements when instructions are jointly used, significantly outperforming all baselines. In fact, when increasing the number of instructions per task from 0, 1 to 8.3, performance of MetaICL\ improves much more than performance of Multi-task 0-shot does. To summarize, (1) learning to in-context learn (MetaICL) outperforms learning to learn from instructions; (2) MetaICL\ and using instructions are largely complementary, and (3) MetaICL\ actually benefits more from using instructions than Multi-task 0-shot does.

Importantly, Channel MetaICL\ trained on available tasks and instructions still achieves lower performance than Channel MetaICL\ without templates/instructions ($46.9$ from Table~\ref{tab:ablate_inst} vs. $49.1$ from Table~\ref{tab:main-result}). This is likely because the model with instructions was trained with less meta-training tasks, which was unavoidable since instructions are only available on 32 out of 61 meta-training tasks. This supports our earlier choice of not using human-written templates/instructions, since writing templates and instructions for every task requires extensive effort.

It is worth noting that, it is nonetheless difficult to make direct comparisons with \citet{wei2022finetuned} and \citet{sanh2022multitask} because there are many moving components: size of LMs, types of LMs (e.g., causal LM vs. masked LM), splits between meta-training and target tasks, and more.

\section{Conclusion}\label{sec:concl}
In this paper, we introduced MetaICL, a new few-shot learning method where an LM is meta-trained to learn to in-context learn, i.e. condition on training examples to recover the task and make predictions.

We experiment with a large, diverse collection of tasks, consisting of 142 unique tasks in total and 52 unique target tasks, using seven different settings.
MetaICL\ outperforms a range of strong baselines including in-context learning without meta-training and multi-task learning followed by zero-shot transfer, and outperforms or matches 8x bigger models. We identify ingredients for success of MetaICL\ such as the number and diversity of meta-training tasks. We also demonstrate that, while MetaICL\ is better than recent work using natural instructions, they are complementary and the best performance is achieved by integrating MetaICL\ with instructions.

\paragraph{Limitation \& Future work} Our work is limited in multiple dimensions. First, in-context learning approaches in general requires much longer context at both meta-training and inference due to feeding the concatenation of the training data, thus being less efficient compared to baselines that do not use in-context learning. Second, our work experiment with a casual language model with modest size (GPT-2 Large, 770M parameters). Future work may investigate extending our approach to a masked language model and a larger model. Third, our experiments focus on classification and multi-choice tasks where a set of candidate options is given. Future work may study applying our approach for a wider range of tasks including free-form generation.

Other avenues for future work include further improving MetaICL\ to outperform supervised models with meta-training, identification of which meta-training tasks are helpful on target tasks, and how to better combine human-written instructions and MetaICL. 

\end{document}
